\documentclass[11pt,a4paper]{article}
\usepackage[margin=25mm]{geometry}
\usepackage{amsmath,amssymb,amsthm}
\newtheorem{lemma}{Lemma}\newtheorem{proposition}[lemma]{Proposition}
\newcommand{\Sz}{\operatorname{Sz}}
\title{Descent of the order-four subfield configuration}
\author{Human Compression project: F003 working continuation}\date{}
\begin{document}\maketitle
This note continues the accepted proper-face, mixed-face and cyclic-overlap
arguments for the fixed F003 source. It is a bounded reduction, not a
recovery of the whole rank theorem. We retain their classical Suzuki
structure and standard subfield interfaces, with the same attribution.
The coordinate descent below is an additional argument. Distinct review
has not yet been completed.

Let $S=\Sz(2^N)$, where $N$ is odd, and put
$\theta=2^{(N+1)/2}$. Thus $(r^\theta)^\theta=r^2$ for every ambient
field element $r$. Let $K=\mathbb F_{2^d}$ and
$L=\mathbb F_{2^e}$ be proper subfields, with $d,e>1$. Suppose
\[
 F=\langle a,b\rangle\cong\Sz(2^d),\qquad
 E=\langle b,c\rangle\cong\Sz(2^e),\qquad F\cap E=\langle b\rangle,
\]
where $ab$, $bc$ and $abc$ are involutions and $|b|=4$.
The accepted cyclic-overlap argument gives $\gcd(d,e)=1$.
We use the standard matrices
\[
 U(x,y)=\begin{pmatrix}
 1&0&0&0\\x&1&0&0\\x^{\theta+1}+y&x^\theta&1&0\\
 x^{\theta+2}+xy+y^\theta&y&x&1
 \end{pmatrix},\quad z(t)=U(0,t),\quad
 J=\begin{pmatrix}0&0&0&1\\0&0&1&0\\0&1&0&0\\1&0&0&0\end{pmatrix}.
\]
Direct multiplication gives
$U(x,y)U(z,w)=U(x+z,y+w+xz^\theta)$.

\begin{lemma}[Common normalization]\label{normal}
After simultaneous conjugation, one may assume
\[
 F=\Sz(K),\qquad b=U(1,\epsilon),\qquad
 E=\Sz(L)^{z(t)},\qquad \epsilon\in\{0,1\},
\]
where $t\notin K+L$. In these coordinates, after changing $t$ by
an element of $K+L$ if necessary, the involutions $u=ab$ and $v=bc$
have the forms
\[
 u=U(x,0)^{-1}Jz(r)JU(x,0),\qquad
 v=U(y,t)^{-1}Jz(s)JU(y,t),
\]
with $x,r\in K$, $y,s\in L$, and $r,s\ne0$.
\end{lemma}
\begin{proof}
The unique fixed point of $b$ belongs to both subovoids by the retained
fixed-point descent lemma. Conjugate $F$ to the standard subgroup and
this point to infinity. Within each subfield the diagonal torus
normalizes the nonzero first coordinate of an order-four element to
one. Conjugation by $U(h,k)$ changes the second coordinate of
$U(1,\eta)$ by $h+h^\theta$. The image of $h\mapsto h+h^\theta$
is exactly the absolute-trace-zero hyperplane: its kernel is
$\mathbb F_2$, and its image has trace zero. Since the field degree
is odd, the two trace classes are represented by $0$ and $1$.
Ambient extension of odd degree preserves their distinction.
Consequently the two standard subfield representatives of $b$ can
both be taken to be the same $U(1,\epsilon)$.

Their remaining conjugator centralizes $b$. It fixes infinity and
therefore belongs to the point stabilizer. Its torus part is trivial,
since it preserves the first coordinate one. Its root-group first
coordinate lies in $\mathbb F_2$, since $h+h^\theta=0$. Hence
\[
 C_S(b)=\{U(h,k):h\in\mathbb F_2,\ k\in\mathbb F_{2^N}\}
       =\langle b\rangle Z(U).
\]
Absorb the power of $b$ into the smaller standard group; its other
factor is $z(t)$. If $t=k+l$ with $k\in K$, $l\in L$, conjugation
by $z(k)$ preserves $F$ and takes $E$ to its standard embedding.
Their intersection then contains $\Sz(2)$, of order twenty,
contrary to its being $\langle b\rangle$ of order four.

Neither $u$ nor $v$ fixes infinity: otherwise its face group fixes
infinity and is not a genuine Suzuki subfield group. The root group
acts regularly on the other ovoid points. Every involution with fixed
point $p(x,h)$ is consequently
$U(x,h)^{-1}Jz(r)JU(x,h)$ for a unique nonzero root-centre parameter
$r$. These parameters belong to the appropriate subfield. Conjugating
everything by a suitable $z(k)$, $k\in K$, puts the second coordinate
of the fixed point of $u$ at zero; it preserves $b$ and $F$ and changes
the parameter of $E$ by $k$. The remaining second coordinate of the
fixed point in the standard $L$-subovoid can be absorbed into $t$,
changing it by an element of $L$ without changing $E$. This gives the
stated forms.
\end{proof}

\begin{proposition}[Descent to the compositum]\label{descent}
In the configuration of Lemma~\ref{normal}, $t\in KL$. In particular,
if $\langle a,b,c\rangle=S$, then $N=de$.
\end{proposition}
\begin{proof}
Write $M=KL=\mathbb F_{2^{de}}$. This field is stable under
$\theta$, and its degree in the ambient field is odd. Since
$abc=ub^{-1}v$ is an involution, and $u,v$ are involutions,
\[
 ub^{-1}v=vbu.                                           \tag{1}
\]
We compute only the $(1,3)$ entry of this identity. Write
$R=r^\theta$, $T=t^\theta$, and $S_0=s^\theta$.
Substituting the matrices in Lemma~\ref{normal} shows that this entry
is
\[
 CT+Dt+Q=0,\qquad Q\in M,                               \tag{2}
\]
where, for either value of $\epsilon$,
\begin{align*}
 C&=RS_0(x+y)+Rs+S_0r,\\
 D&=RS_0(x^2+x+y^2+y)+Rs+S_0r.
\end{align*}
In particular there are no quadratic or mixed terms involving $t,T$
in this entry. These coefficient identities follow by multiplying the
displayed four by four matrices. They may also be checked before any
field specialization, in the polynomial ring of characteristic two
with $x,x^\theta,y,y^\theta,r,R,s,S_0,t,T$ treated as independent
symbols. All terms not involving $t,T$ belong to $M$, which is the
only property of $Q$ used here.

The two coefficients cannot both vanish. Indeed, their sum is
$RS_0(x+y)^2$. If they both vanished, then $x=y$, and $C=0$ would
give $r/R=s/S_0$. The first common value belongs to
$K\cap L=\mathbb F_2$, and the second belongs to
$K^\times\cap L^\times=\{1\}$. Thus $x\in\mathbb F_2$ and
$r=r^\theta$. The fixed field of $\theta$ on $K$ is
$\mathbb F_2$: the ambient twist restricts to
$2^{(d+1)/2}$ and $\gcd(d,(d+1)/2)=1$. Hence $r=1$.
The displayed matrix for $u$ then has entries in $\mathbb F_2$,
as does $b$. This places $F=\langle u,b\rangle$ inside the standard
$\Sz(2)$, contradicting $d>1$.

If $C=0$, equation (2) immediately gives $t\in M$.
Otherwise write (2) as $t^\theta=\alpha t+\beta$ with
$\alpha,\beta\in M$. Applying $\theta$ gives
\[
 t^2=\alpha^\theta\alpha t+
       \alpha^\theta\beta+\beta^\theta.
\]
Thus $t$ has degree at most two over $M$. Its degree also divides the
odd degree $[\mathbb F_{2^N}:M]$, so it is one. This proves
$t\in M$ in the remaining case.

All the matrices $u,b,v$ now belong to the standard $\Sz(M)$.
Since $a=ub^{-1}$ and $c=b^{-1}v$, so do the original three
generators. If they generate $S$, the subfield cannot be proper;
therefore $N=de$.
\end{proof}

This leaves the case $N=de$ and $t\in KL\setminus(K+L)$ unresolved.
The case of an odd-order middle rotation also remains unresolved. No conclusion
about those cases follows from this proposition, and no whole-original,
conjecture, formalization or compression credit is claimed.
\end{document}
